% New module; existing joint-continuation and primitive appendices are unchanged.
\subsection{From an accepted ContextBound operation to the message continuation}
\label{sec:contextbound-connection}

We now discharge the ContextBound comparison left explicit in
Section~\ref{sec:joint-continuation} for one concrete native profile.
The protected component is ML-KEM-768 with a real-versus-uniform 32-byte
shared-secret challenge. The other component is evaluated by the actual
X25519 algorithm. Its raw DH output is not assumed to satisfy that same
uniform-key experiment. Write $H$ for the whole SHA3-256 random oracle and
$h$ for the independent SHA-256 compression oracle of
Section~\ref{sec:joint-continuation}. Both remain available throughout every
comparison. All oracle-dependent preparation, including construction of an
established starting state, is executed and counted; an oracle-correlated
history is not supplied as free auxiliary input.

\paragraph{Complete effective tuples.}
For a byte string $x$, let
$\operatorname{LP}(x)=\operatorname{BE64}(|x|)\Vert x$.
The admitted native ContextBound computation is $S=H(\operatorname{enc}(T))$,
where $\operatorname{enc}$ concatenates the LP encodings of the ten fields
\[
 T=(D_{\rm CB},\mathit{suite},\operatorname{BE32}(v),
 s_M,s_X,c_M,p_M,c_X,p_X,C_{\rm policy}).
\]
Here
\[
\begin{aligned}
 C_{\rm policy}
 &=\operatorname{LP}(D_{\rm policy})
   \Vert\operatorname{LP}(\mathit{policyDigest})
   \Vert\operatorname{LP}(C_{\rm app}),\\
 C_{\rm app}
 &=H\bigl(\operatorname{LP}(D\Vert\texttt{response-kem})
                \Vert\operatorname{LP}(O)\bigr),
\end{aligned}
\]
with $O$ the entire retained signed offer and the policy digest that of the
verified admitted policy. In the pinned suite the complete encoded input
has 2639 bytes, and the 32-byte $s_M$ field occupies offsets 76--107.
The reservation's trust-root binding, original-slot identifiers, operation
identifiers and journal identities are separate metadata; they are not
extra fields of this hash input. In particular, the raw offer is present
through $C_{\rm app}$, not additionally appended to the outer tuple.

\begin{lemma}[Complete-tuple hidden-component comparison]
\label{lem:cb-hidden-component}
Let $U$ be uniform on $\{0,1\}^{256}$, independent of the initially sampled
oracles and initial auxiliary state. In a protected family replace one
32-byte component field by $U$, and let $m$ contain all nine other actual
fields, including the other component secret and the complete policy
context. The private interface returns
$H(f(U,m))$, where $f$ is the actual ten-field encoding. Exact repeats are
cached by $m$. Assume that the rest of the experiment accesses $U$ only
through this interface or opaque handles and respects its stated stops.
For at most $d$ distinct private inputs and $Q_H$ ordinary queries to $H$,
replacing the private interface by a random function $\mathcal T(m)$
changes every efficient output-bit probability by at most
\begin{equation}
 \varepsilon_{\rm CB}=\Phi_{256}(d,Q_H).
 \label{eq:cb-hidden-cost}
\end{equation}
The inputs $m$ may be adaptive. $Q_H$ counts the complete derived
experiment, including preparation, all ordinary KEM/hash work, continuation
and terminal test. The same independent oracle $h$ is retained.
\end{lemma}
\begin{proof}
Length-prefix parsing recovers all ten actual fields. For fixed $m$ and any
complete string $w$, at most one $u\in\{0,1\}^{256}$ satisfies
$f(u,m)=w$, so its point probability is at most $2^{-256}$. If $m\ne m'$
as effective field tuples, then $f(u,m)\ne f(u,m')$ for every $u$.
Cache repeats before invoking the private interface. These are precisely
the point-probability and all-secret freshness conditions of the
hidden-input theorem stated in Appendix~\ref{app:primitive-inputs}
\cite[Theorem~2]{dfh2022}. It gives the displayed bound for adaptive
private queries. The independent $h$ interface is part of the unchanged
auxiliary computation. Ordinary $H$ queries, including an accidentally
matching known computation, retain their actual oracle responses and are
counted in $Q_H$. The argument does not condition on avoiding a collision
or on the final acceptance event.
\end{proof}

The cache key in this lemma cannot be replaced by an origin identifier.
Two different original calls can have the same effective tuple: their real
answers agree with probability one, whereas two independent per-origin
ideal answers agree with probability $2^{-256}$. Conversely, dropping the
other shared secret or any effective field could merge distinct real hash
inputs. Provenance and hash consistency therefore require separate tables.

\paragraph{Accepted original component calls.}
Fix $j=(g,a)$ in the existing original-reservation ledger. The component
adapter assigns its ML-KEM challenge public key only to generation origin
$g$, retaining the locally sampled X25519 scalar, runtime admission and
owner metadata. It obtains its challenge pair $(c_M^*,K_{\rm ch})$ at the
first actual private encapsulation computation of $a$ on the original
retained offer from $g$. It computes the X25519 leg locally and caches the
whole result before subsequent persistence. Exact retries reuse that
challenge; a fresh entropy reservation, including one after an unsuccessful
attempt, receives its actual new original slot. Copies and reopenings do
not create a new challenge invocation.

After challenge creation, a PQ decapsulation under the selected generation
owner with ciphertext $c_M^*$ returns the cached $K_{\rm ch}$, regardless
of accompanying traditional ciphertext or context. All other PQ
ciphertexts are routed to the component decapsulation interface. The
traditional decapsulation and its rejection path remain native; a rejected
traditional leg produces no ContextBound call. The challenged PQ
ciphertext is never submitted to the postchallenge CCA oracle. The operational component game first supplies its generated public key
and permits an ordinary prechallenge decapsulation phase. A single
challenge is requested at the specified first computation of $a$; it
returns a fresh encapsulation ciphertext and either its real key or an
independent uniform 32-byte key. Thereafter its decapsulation oracle
rejects that challenge ciphertext. Before challenge creation, selected-key
decapsulation follows this ordinary prechallenge interface; a query that
happens to equal the later sampled $c_M^*$ retains its original answer
and is not deleted or retroactively replaced by $U$. Other original
owners use their locally known keys, even when public-key bytes happen
to coincide. No equivalence with an interface that supplies its challenge
up front is assumed without a separate game conversion.

The immutable admitted binding must show that $a$ consumes the original
offer from $g$. An immutable contrary binding makes attribution to $j$
permanently impossible and gives a free stop $U_j$ under
Section~\ref{sec:candidate-local-accounting}; byte equality alone does not
transfer provenance. Missing or unaccepted operations otherwise retain the
native execution. Reveal and projection cases use the established free
stops, charged stops and local record correspondence. Thus the existing
component costs $a_j^+,a_j^-,\delta_j$ retain exactly their meanings;
this adapter does not set correctness or permitted-projection costs to zero.
More explicitly, the derived distinguisher $D_j^\sigma$ runs the entire
adapter and outputs $\mathbf1[\Psi_j=\sigma]$ for
$\sigma\in\{-1,+1\}$. Its operational advantage is
\[
 a_j^\sigma=\left|\Pr[D_j^\sigma=1\mid\text{real challenge}]
                    -\Pr[D_j^\sigma=1\mid\text{uniform challenge}]\right|.
\]
The real branch matches $\mathsf V^{[j]}$ up to the selected-pair
correctness cost $\delta_j$; the uniform branch is exactly the declared
$\widetilde{\mathsf V}^{[j]}$. This gives the two output-event
inequalities used in Proposition~\ref{prop:native-local-accounting}.

In the resulting random-component game,
$\widetilde{\mathsf V}^{[j]}$, $K_{\rm ch}=U$ is independent uniform.
The reduction for Lemma~\ref{lem:cb-hidden-component} can materialize all
actual native KEM tapes locally, execute their algorithms, and replace
only the designated protected component returns by $U$. This is a new
implementation of that adjacent pair of experiments, not reuse of the
previous reduction's unavailable challenge tapes. Ordinary SHA3 work
inside those algorithms is included in $Q_H$; opaque records and the same
stopping code retain their declared interfaces.

\paragraph{A concrete single-root source profile.}
Use the next rekey of one established authoritative session. The two
endpoints initially have the same established root $r$, confirmed epoch
and control parent, with no pending selected rekey. The old root and
wrapping keys may be known. The supported native call surface includes
failed reservations and writes, process loss and exact resumption,
adversarial public packets, old-epoch traffic, asynchronous final/receipt
processing, the marked first-message interface of
Corollary~\ref{cor:joint-first-message}, and subsequent epochs. Private
selected owners are not additionally exported to arbitrary low-level SDK
callers, and read-only copies are not activated as competing histories.
These are an initial state and call surface; the original eligibility
predicate is unchanged.

\begin{lemma}[Next-rekey root-use correspondence]
\label{lem:cb-native-root-use}
In this profile, fix $j=(g,a)$ whose immutable original encapsulation
binding consumes the retained original offer from $g$ for this next rekey.
Let $m^*$ be that original encapsulation's effective tuple: its challenged
$c_M^*$, locally generated $c_X^*$, traditional secret, retained offer and
admitted policy determine all nine fields. Every protected
use of that same tuple feeds one pair $(r,I_R)$, where $I_R$ is the complete
97-byte native pending-root information string. They are one original
Extract/Expand and exact repetitions. No raw combined-secret channel is
added by retry, crash or completion. This invariant holds on every
continuation of the matching binding, independently of later eligibility
or attribution.
\end{lemma}
\begin{proof}
The retained generation token proceeds to a fixed signed offer. Its
reconstruction checks the regenerated public material against the saved
offer. At the responder the encapsulation token is bound to the complete
saved offer; a different offer conflicts, and completing or committed
plans return their retained flights. While these tokens are pending,
ordinary message operations do not change the established rekey root.
The response prefix fixes the session, context, epochs, role, control
parent and retained-offer digest. Equality of $m^*$ fixes both complete
ciphertext fields, hence the response core and $I_R$. Both endpoints
therefore use their common old $r$ and this same $I_R$.

Before exposing its response, the responder replaces the active
encapsulation-token stage by the stage retaining the derived root.
The proposer replaces the generation-token stage before root installation.
The final/receipt installation sites are the live updates of the rekey
root. A crash before a transition recovers the same token and old root; a
committed transition recovers the retained result. Neither resumes an old
selected token under a newly installed root. This invariant is preserved
by every supported call, authoritative write outcome and recovery step.
The source artifact records the corresponding pinned call sites and the
complete induction. The statement concerns protected-family uses; it
makes no claim that independently sampled numeric secrets never coincide.
\end{proof}

\paragraph{Fresh allocation and score semantics.}
Mark this original target tuple $m^*$ before its first actual private hash
computation, at any protected alias call site and including a computation
whose subsequent write fails. If a crash intervenes after component
challenge creation, its cached ciphertext and locally determined traditional
leg still identify $m^*$ before any such first use. An earlier private use
at any alias site already allocates this entry and cannot be ignored at a
later durable acceptance or network release. Other protected tuples do not
become the target. If the selected original $a$ never supplies a usable
target, the original candidate score is zero; no substitute tuple is chosen.
In the private-random-function game,
$S^*=\mathcal T(m^*)$ is a freshly allocated uniform 32-byte value;
replays and aliases retrieve that same entry. It can be presampled and
withheld at this allocation. This is an online allocation argument,
not a claim of uniformity after conditioning on the future transcript.
Other entries of $\mathcal T$ can be sampled by the controller. By
Lemma~\ref{lem:cb-native-root-use}, all selected uses of $S^*$ are the
one Extract/Expand required by Appendix~\ref{app:joint-root-entry}.
The actual continuation is then the controller of
Section~\ref{sec:joint-continuation}, including its fresh traffic-entry
and direction mark, confirmations, ACKs, successors, shared $h$ work and
private retained state.

Every hop runs the same efficient terminal code $\Psi_j$ on its own actual
trace. Its value is zero after the established local stops. The qualifying
predicate in the message game must be computable from the available
classical view and provenance ledger, as required by
Corollary~\ref{cor:joint-first-message}. A hidden challenge bit is not an
additional input to that predicate. This condition matters: if the guess
is always zero and a fictitious predicate accepts exactly when the hidden
bit is zero, the qualified signed score has expectation $1/2$ even for an
otherwise bit-independent view. For the admitted predicate, making an
independent fair guess on disqualified runs implements the full qualified
score in the actual AEAD game; no unsupported zero-mean idealization is
needed.

\begin{theorem}[ContextBound connection for the next native rekey]
\label{thm:cb-next-rekey}
Use the concrete profile and adapters above, the local comparisons of
Proposition~\ref{prop:native-local-accounting}, and the message interface
of Corollary~\ref{cor:joint-first-message}. Fix finite budgets before the
experiment. For candidate $j$, let $d_j$ bound its distinct protected
effective tuples and $Q_{H,j}$ its complete ordinary $H$ queries. Let
$\Delta_{{\rm init},j}$, $\Delta_{{\rm sched},j}$ and
$\varepsilon_{{\rm AEAD},j}$ have the existing whole-continuation meanings
of Appendix~\ref{app:joint-root-entry},
Theorem~\ref{thm:joint-message-leaf} and
Corollary~\ref{cor:joint-first-message}. Define
\begin{equation}
 B_j^{\rm cont}=\min\left\{1,
 2\Phi_{256}(d_j,Q_{H,j})+2\Delta_{{\rm init},j}
 +2\Delta_{{\rm sched},j}+2\varepsilon_{{\rm AEAD},j}\right\}.
 \label{eq:cb-continuation-bound}
\end{equation}
Then $|b_j|\le B_j^{\rm cont}$ and
\begin{equation}
\begin{split}
 \operatorname{Adv}_{\rm msg}\le
 \min\!\left\{\frac12,\frac12\left(
 \eta_{\rm attr}^{\rm nat}+\eta_{\rm proj}^{\rm nat}\right.\right.\\[-2pt]
 \left.\left.\qquad+
 \sum_{j\in J}\bigl(\rho_j+a_j^++a_j^-+2\delta_j+B_j^{\rm cont}\bigr)
 \right)\right\}.
\end{split}
 \label{eq:cb-native-bound}
\end{equation}
The residuals are still evaluated in the one original native experiment.
Each candidate's ideal continuation is compared separately.
\end{theorem}
\begin{proof}
Apply Lemma~\ref{lem:cb-hidden-component} to
$\widetilde{\mathsf V}^{[j]}$. Fresh allocation of $m^*$ gives precisely
the uniform-combined-contribution game
$\mathsf W_{\rm cont}^{[j]}$ from
Section~\ref{sec:joint-continuation}; all other private table entries and
the entire native controller remain present. For any
$\Psi_j\in[-1,1]$, a terminal randomized test that outputs one with
probability $(\Psi_j+1)/2$ converts a bit gap $\varepsilon$ into a signed
score gap at most $2\varepsilon$. Thus
\[
 |b_j-b_j^{\rm cont}|\le2\Phi_{256}(d_j,Q_{H,j}).
\]
The native root-use lemma supplies the single-entry hypothesis for the
existing $\Delta_{{\rm init},j}$ comparison. The joint schedule and AEAD
continuation then give
$|b_j^{\rm cont}|\le2\Delta_{{\rm init},j}
 +2\Delta_{{\rm sched},j}+2\varepsilon_{{\rm AEAD},j}$.
All derived distinguishers include the complete preparation and continuation;
no oracle budget resets at a hop. Combine these inequalities, use
$|b_j|\le1$, and substitute in
Proposition~\ref{prop:native-local-accounting}.
\end{proof}

The result supplies the previously explicit ContextBound comparison for the
next-rekey instance. Its source requirement is expressed through the root-use
descriptor, so extensions can retain the same tuple-level comparison while
establishing their corresponding root-entry interface.
